% paper/tables/inverse.tex -- generated by analyze/paper/tables_paper.py; sources in paper/tables/README.md
\setlength{\tabcolsep}{3.5pt}%
\begin{tabular}{@{}lllll@{}}
\toprule
\multicolumn{5}{@{}p{0.9\linewidth}@{}}{(a) 自己視界基準の正解率 BA（$T=50$，L1 期間の最終5点の範囲，括弧は合格の点）} \\
\midrule
組 & 経験80\% & 経験50\% & 経験なし & \\
\midrule
組1 & 0.641〜0.671（5/5） & 0.598〜0.603（5/5） & 0.389〜0.420 & \\
組2 & 0.667〜0.701（5/5） & 0.526〜0.546（0/5） & 0.499〜0.518 & \\
組3 & 0.553〜0.593（4/5） & 0.445〜0.483（0/5） & 0.506〜0.517 & \\
組4 & 0.532〜0.564（2/5） & 0.486〜0.507（0/5） & 0.469〜0.494 & \\
組5 & 0.589〜0.618（0/5） & 0.534〜0.567（0/5） & 0.557〜0.574 & \\
組6 & 0.626〜0.653（5/5） & 0.479〜0.513（0/5） & 0.392〜0.432 & \\
\midrule
方角だけの基準 & \multicolumn{4}{l}{0.708 [0.653, 0.759]（素の一致率 0.707）} \\
 & \multicolumn{4}{l}{規則の一歩が少ない3分の1：0.614 [0.522, 0.704]（117エピソード，素の一致率 0.607）} \\
\midrule[\heavyrulewidth]
\multicolumn{5}{@{}p{0.9\linewidth}@{}}{(b) 経験80\%条件の他者特徴基準の合格の点，好み追従率（他者特徴基準は MSE 期間，自己視界基準は L1 期間の最終5点の範囲，括弧は CI の下限 $>0.5$ の点）} \\
\midrule
組 & 他者特徴基準 & 好み追従率 & 好み追従率 & \\
 & 合格の点 & 他者特徴基準 & 自己視界基準 & \\
\midrule
組1 & 5/5 & 0.489〜0.580（0/5，満たさなかった） & 0.602〜0.670（4/5） & \\
組2 & 5/5 & 0.670〜0.727（5/5，満たした） & 0.591〜0.670（3/5） & \\
組3 & 4/5 & 0.602〜0.659（4/5，満たした） & 0.523〜0.557（0/5） & \\
組4 & 5/5 & 0.602〜0.659（3/5，満たした） & 0.511〜0.523（0/5） & \\
組5 & 5/5 & 0.580〜0.648（2/5，満たさなかった） & 0.500〜0.545（0/5） & \\
組6 & 5/5 & 0.568〜0.625（3/5，満たした） & 0.557〜0.591（0/5） & \\
\midrule[\heavyrulewidth]
\multicolumn{5}{@{}p{0.9\linewidth}@{}}{(c) 経験80\% $-$ 非共有（自己視界基準の BA，$T=50$．組1〜3）：合格の点，ep200 の差 [95\%信頼区間]} \\
\midrule
組1 & \multicolumn{4}{l}{5/5，$+0.255$ [$+0.185$, $+0.324$]} \\
組2 & \multicolumn{4}{l}{5/5，$+0.157$ [$+0.088$, $+0.229$]} \\
組3 & \multicolumn{4}{l}{5/5，$+0.107$ [$+0.037$, $+0.177$]} \\
\bottomrule
\end{tabular}
